%
\section{Factorization of Commutatively zero polynomials}\label{cz-sec}

In this section we will describe the general case of the factorization
algorithm when the input polynomial $f\in\fR$ is a commutatively zero
polynomial. That is, $f$ evaluates to zero on all scalar substitutions
from $\F_q$ or any (commutative) extension field.

%
%
%

The factorization algorithm will follow the three broad steps
described in Section \ref{cnz-sec} for the commutatively nonzero case:
first, using Higman linearization and Theorem \ref{full-monic},
transform the polynomial $f$ to a stably associated linear matrix $L$
that is full and left (or right) monic. Next, factorize the linear
matrix $L$ into atoms. Finally, recover the irreducible factors of $f$
from the atomic factors of the linear matrix $L$ using the factor
extraction procedure described in Lemma~\ref{extract}.

The step that requires a new algorithm is factorizing a full and right (or left)
monic linear matrix $L\in\fR$ into atoms when $f$ is commutatively
zero, which means there is no scalar substitution $x_i\leftarrow
\alpha_i, i\in[n]$ such that $L(\alpha_1,\alpha_2,\ldots,\alpha_n)$ is
invertible. Note that in this case we cannot apply the algorithm for
factorizing a linear matrix as discussed in the proof of Theorem
\ref{lfact1}).

\subsection{Factorization of full and monic linear matrices}

Let $f\in \fR$ be the input polynomial given by a size $s$ formula and
let $L \in\fR^{d\times d}$ be a full, right monic linear matrix stably
associated with $f$ obtained via Higman linearization and an
application of Theorem \ref{full-monic}.

Recall, by Equation~\ref{eq2} we have $f \oplus I_s= PU(L \oplus
I_t)SQ$ where, $P$, $Q$ are respectively upper triangular and lower
triangular units with diagonal entries $1$, $U$ is a unit and $S$ is
scalar invertible. 

%
%
%
%
%
%

Let $L=A_0+\sum_{i=1}^n A_n x_i\in\fR^{d\times d}$ be the given full
and right monic linear matrix.  First, by Lemma~\ref{shift-inv-lemma},
we will find a suitable scalar matrix $n$-tuple
$\bar{M}=(M_1,M_2,\ldots,M_n)$, each $M_i\in \F_q^{\ell\times \ell}$
for $\ell\le 2d$, such that under the substitution $x_i\leftarrow M_i$
the matrix $L(\bar{M})$ is invertible.

For $1\le i\le n$ let $Y_i$ be an $\ell\times \ell$ matrix of distinct
noncommuting variables $y_{ijk}$. We consider the dilated linear
matrix
\begin{equation}\label{dilate-eq1}
L' = A_0\otimes I_\ell + \sum_{i=1}^n A_i\otimes (Y_i+M_i).
\end{equation}
It is not hard to see that $L'$ is full and $L'$ is right monic as  $L$ is right monic. Additionally, its constant term is 
invertible. So, we can apply Theorem~\ref{lfact1} to factorize $L'$ as
a product of two linear matrices, both non-units.

The following lemma \cite{HKV20} has an important role in our
algorithm for recovering the factorization for $L$ from a
factorization of $L'$.

\begin{lemma}\cite{HKV20}\label{hkv20-main-lemma}

Let $L \in \fR^{d \times d}$ be a full linear matrix with $L= A_0+ A_1x_1+\ldots +A_n x_n$ such that $A_i \neq 0$ for at least one $i$, $1\leq i \leq n$ and $L'\in
R^{d\ell \times d\ell}$ be a matrix obtained from $L$ by substituting
variable $x_i$ by $Y_i$ for $i \in [n]$, where $Y_i$ is $\ell \times
\ell$ matrix whose $(j,k)^{th}$ entry is a fresh noncommuting variable
$y_{i,j,k}$ for $1 \leq j,k \leq \ell$. Then
\begin{enumerate}
\item If $L'$ is of the form $G L' H = \left(
\begin{array}{c|c}
A' & 0 \\
\hline
D' & B'
\end{array}
\right)$, where $A'$ is $d' \times d'$ matrix and $B'$ is $d'' \times d''$ matrix for $0< d', d''$, with $d'+d''=d\ell$ and $G,H$ are $d\ell \times d\ell$ invertible scalar matrices then there exist $d \times d$ invertible scalar matrices $U, V$ such that  $U L V = \left(
\begin{array}{c|c}
A & 0 \\
\hline
D & B
\end{array}
\right)$, where $A$ is $e' \times e'$ matrix and $B$ is $e'' \times e''$ matrix for $0<e', e''$, with $e' +e''= d$. 

\item Moreover, given $L'$ explicitly along with its representation
  mentioned above, we can find the matrices $U, V$ in deterministic
  polynomial time (in $n,\ell,d$).
\end{enumerate}
\end{lemma}

\begin{remark}
We give a self-contained complete proof of the above linear-algebraic
lemma in the appendix for $\F_q$, because the proof given in
\cite{HKV20} is sketchy in parts with some details missing, and also
their lemma is stated only for complex numbers and they are not
concerned about computing the matrices $U$ and $V$.
\end{remark}

Now, we can apply Lemma \ref{hkv20-main-lemma} to transform the
factorization of $L'$ to a factorization of $L$ as a product of two
linear matrices, both non-units. Repeating the above on both the
factors of $L$ we will get a complete atomic factorization of
$L$. Formally, we prove the following.

\begin{theorem}\label{lfact2}
On input a full and right (or left) monic linear matrix $L= A_0 +
\sum_{i=1}^n A_i x_i$ where $A_i \in \mathbb{F}^{d \times d}$ for $i\in
    [n]$, there is a randomized polynomial time ($poly(n,d)$)
    algorithm to compute scalar invertible matrices $S,S'$ such that
    $SLS'$ has atomic block diagonal form.
\end{theorem}

\begin{proof}
  We present the algorithm only for right monic $L$; the left monic
  case has essentially the same solution.

  If the input $L$ is not full or right monic the algorithm can
  efficiently detect that and output ``failure''. If $L$ is an atom
  the algorithm will output that $L$ is an atom and set the matrices
  $S$ and $S'$ to $I_d$. Otherwise, the algorithm will compute
  invertible scalar matrices $S$ and $S'$ such that
\begin{equation}\label{eq-block-atom}
S L S' = 
\left(
\begin{array}{ccccc}
L_1 & 0 &0 &\ldots & 0\\
* & L_2 &0 &\ldots & 0\\
* & * &L_3 &\ldots & 0\\
& & &\ddots &\\
* & * &* &\ldots & L_r\\
\end{array}
\right),
\end{equation}
where the matrix on the right is in atomic block diagonal form, that
is, each linear matrix $L_i$ is an atom.

\begin{enumerate}
\item[]   {\bf Procedure Factor(L)}.
\item Test if $L$ has full noncommutative rank using the algorithm in
  \cite{IQS17} or \cite{GGOW20}. Test if $L$ is right monic by checking if
  the matrix $[A_1 A_2 \ldots A_n]$ has full row rank (which is $d$).
  If $L$ is not full and right monic the algorithm outputs ``fail''.
\item Assume $L$ is full and right monic. Using Lemma
  \ref{shift-inv-lemma}, find smallest positive integer $\ell \leq 2d$
  and $\ell\times \ell$ scalar matrices $M_i, i\in [n]$ with entries
  from $\F$ (or a small degree extension of $\F$) such that
  $W=L(\bar{M})$ is $d \cdot \ell \times d \cdot \ell$ invertible
  scalar matrix. Compute the dilated linear matrix $L'$ in the
  $y_{ijk}$ variables as in Equation~\ref{dilate-eq1} which can be
  rewritten as:
\[
L' = A_0\otimes I_\ell + \sum_{i=1}^n A_i\otimes M_i +
\sum_{i=1}^n\sum_{j,k=1}^\ell (A_i\otimes E_{jk})\cdot y_{ijk}.
\]
Let $L''= W^{-1} L'$. Clearly $L''(\overline{0})= I_{d\ell}$. Hence,
by the algorithm of Theorem \ref{lfact1} we can either detect that
$L''$ is an atom or factorize $L''$. If $L''$ is an atom then $L$ is
also an atom and the algorithm can output that and stop.  Otherwise,
$L'$ is not an atom and by Theorem \ref{lfact1} we will obtain a basis  
change matrix $T$ such that $T W^{-1} L' T^{-1}= T L'' T^{-1} = \left(
\begin{array}{cc}
C'' & 0 \\
* & D''
\end{array}
\right)$ where $C''$ and $D''$ are linear matrices of dimension $c''
\times c''$ and $d'' \times d''$ respectively, such that $c''+ d'' = d
\ell$.

%
%
%
%
%
%
  
\item By linear shift of variables $y_{ijk} \leftarrow y_{ijk} -
  M_i(j,k)$ we obtain $\tilde{T} \tilde{L} \tilde{T'} = \left(
\begin{array}{cc}
C' & 0 \\
* & D'
\end{array}
\right)$    
for some scalar invertible matrices $\tilde{T}, \tilde{T'}$ where
$\tilde{L} = L( Y_1, \ldots, Y_n)$.

\item Applying the algorithm of Lemma \ref{hkv20-main-lemma} to
  $\tilde{L}$, $\tilde{T}$, and $\tilde{T'}$, in deterministic
  polynomial time we obtain scalar invertible matrices $\tilde{S}, \tilde{S'}$ such
  that $\tilde{S}L\tilde{S'} = \left(
\begin{array}{cc}
C & 0 \\
* & D
\end{array}
\right) $
where $C$, $D$ are square matrices of dimensions $e \times e$ and $g
\times g$, respectively, such that $e+g = d$.

\item Recursively call Factor$(C)$ and Factor$(D)$. Let $S_1, S'_1$ be
  the matrices returned by Factor$(C)$ and $S_2, S'_2$ be the matrices
  returned by Factor$(D)$.

\item Let $S =(S_1 \oplus S_2)\tilde{S}$ and $S' = \tilde{S'}(S'_1
  \oplus S'_2)$. Return the invertible scalar matrices $S$ and
  $S'$. Note that at this stage $SLS'$ has the desired atomic block
  diagonal form.
\end{enumerate}

Next we give a brief argument for proving correctness of the above
algorithm. Firstly, the algorithm declares $L$ as an atom iff $L$ is
indeed an atom. To see this, we will prove $L$ is not an atom iff
$L''$ is not an atom. Forward direction is obvious. To prove the
reverse direction of implication, let $L''$ is not an atom. Which
implies $L' = W L''$ is not an atom. $\tilde{L}$ is a linear matrix
obtained by substituting $M_i(j,k)=0$ for all $i,j,k$ in
$L'$. Clearly, $\tilde{L}$ is not an atom as $L'$ is not an
atom. Using Lemma \ref{hkv20-main-lemma} it follows that $L$ is not an
atom. So we have established $L$ is not a atom iff $L''$ is not an
atom. So if input linear matrix $L$ is an atom, the algorithm will
correctly declare it to be an atom in step 2.

Now we argue that we will get correct atomic block diagonal form in
the last step of the algorithm. Firstly, for giving recursive calls to
the Factor procedure for the matrices $C$, $D$, we must have $C, D$ to
be right monic as stated in the claim below. This is proved by the
same argument as in the proof of Theorem~\ref{cohnthm}.

%
%

\begin{claim} \label{monic-monic}
Let $L \in \fR^{d \times d}$ be a full and right monic linear matrix
such that $P'LQ' = \left(
\begin{array}{cc}
C & 0 \\
E & D
\end{array}
\right)$ where $C$ and $D$ are linear matrices of dimensions $e \times
e$, $g \times g$, respectively, such that $e+g = d$. Then both $C, D$
are right monic.
\end{claim}

\begin{comment}
\begin{proof}
Let $L=A_0 + \sum_{i=1}^n A_i x_i$ where $A_i \in \mathbb{F}^{d \times
  d}$ for each $i$. Similarly, let $C=C_0 + \sum_{i=1}^n C_i x_i$,
$D=D_0 + \sum_{i=1}^n D_i x_i$, and $E=E_0 + \sum_{i=1}^n E_i x_i$
where $C_i \in \mathbb{F}^{e \times e}$, $D_i \in \mathbb{F}^{g \times
  g}$ and $E_i \in \mathbb{F}^{g \times e}$ for each $i$. As $L$ is
right monic, the matrix $[A_1 A_2 \ldots A_n]$ has full row rank which
implies $[A'_1 A'_2 \ldots A'_n]$ has full row rank, where
$A'_i=PA_iQ$ for each $i$. Since $PA_iQ = A'_i =
\left(
\begin{array}{cc}
C_i & 0 \\
E_i & D_i
\end{array}
\right)$, it follows that $[C_1 0 C_2 0 \ldots C_n ~0]$ has full row
rank, where each $0$ denotes an $e \times g$ sized block of zeros. So
$[C_1~C_2~\ldots ~C_n]$ is full row rank and hence $C$ monic. Now by
elementary row operations applied to matrix $[A_1 A_2 \ldots A_n]$, we
can use the full row rank matrix $[C_1~C_2~\ldots ~C_n]$ (occupying a
subset of columns in the top $e$ rows) to drive the matrix $[E_1 E_2
  \ldots E_n]$ in the corresponding columns of the bottom $g$ rows to
zero without changing the $D_i$ matrices. That means, for an
invertible scalar matrix $R$ corresponding to row operations and
suitable column permutation matrix $S$ we have $R [A_1 A_2 \ldots
A_n] S = \left(
\begin{array}{c|c}
C_1~C_2~\ldots ~ C_n  & 0 \\
\hline
~~~~~~0~~~~~  & D_1~D_2~\ldots ~ D_n
\end{array}
\right)$ Hence $[D_1 D_2 \ldots D_n]$ has full row rank implying that
$D$ is right monic. Similarly, using left monicity of $L$ we can prove
that $C$ and $D$ are left monic. 
\end{proof}
\end{comment}

By recursive calls Factor($C$) and Factor($D$) obtain matrices $S_1,
S'_1, S_2, S'_2$ such that $S_1C S'_1=C'$ and $S_2 D S'_2=D'$ are in
atomic block diagonal form.  We can write $\tilde{S}L\tilde{S'}$ as
\begin{eqnarray*}
&=& \left(
\begin{array}{cc}
C & 0 \\
E & D
\end{array}
\right)\\
&=& \left(
\begin{array}{cc}
C & 0 \\
0 & I_g
\end{array}
\right)\left(
\begin{array}{cc}
I_e & 0 \\
E & I_g
\end{array}
\right)\left(
\begin{array}{cc}
I_e & 0 \\
0 & B
\end{array}
\right)\\
&=& (S_1^{-1} \oplus I_g)(C' \oplus I_g)({S'}_1^{-1} \oplus I_g)\left(
\begin{array}{cc}
I_e & 0 \\
E & I_g
\end{array}
\right)(I_e \oplus S_2^{-1})(I_e \oplus D')(I_e \oplus {S'}_2^{-1})\\
&=&  (S_1^{-1} \oplus I_g)(C' \oplus I_g)(I_e \oplus {S'}_2^{-1})\left(
\begin{array}{cc}
I_e & 0 \\
S_2ES'_1 & I_g
\end{array}
\right)({S'}_1^{-1}\oplus I_g)(I_e \oplus D')(I_e \oplus {S'}_2^{-1})\\
&=& (S_1^{-1} \oplus I_g)(I_e \oplus S_2^{-1})(C' \oplus I_g)\left(
\begin{array}{cc}
I_e & 0 \\
S_2ES'_1 & I_g
\end{array}
\right)(I_e \oplus D')({S'}_1^{-1}\oplus I_g)(I_e \oplus {S'}_2^{-1})\\
&=& (S_1^{-1} \oplus I_g)(I_e \oplus S_2^{-1})\left(
\begin{array}{cc}
C' & 0 \\
S_2ES'_1 & D'
\end{array}
\right)({S'}_1^{-1}\oplus I_g)(I_e \oplus {S'}_2^{-1})\\
&=& (S_1^{-1} \oplus S_2^{-1})\left(
\begin{array}{cc}
C' & 0 \\
S_2ES'_1 & D'
\end{array}
\right)({S'}_1^{-1}\oplus {S'}_2^{-1}).
\end{eqnarray*}
Thus we have
\[
(S_1 \oplus S_2) \tilde{S}L\tilde{S'}(S'_1 \oplus S'_2) = \left(
\begin{array}{cc}
C' & 0 \\
S_2ES'_1 & D'
\end{array}
\right).
\]

As $C'$ and $D'$ are in atomic block diagonal form, it follows that
$\left(
\begin{array}{cc}
C' & 0 \\
S_2ES'_1 & D'
\end{array}
\right)$ is also in atomic block diagonal form. Letting $S=(S_1 \oplus
S_2) \tilde{S}$ and $S'=\tilde{S'}(S'_1 \oplus S'_2)$, it follows that
$SLS'$ is in the desired atomic block diagonal form which proves the
correctness of Factor procedure. In each call to the procedure
(excluding the recursive calls) the algorithm takes
$\poly(n,d,\log_2q)$ time. The total number of recursive calls overall
is bounded by $d$. Hence, the overall running time is
$\poly(n,d,\log_2q)$. This completes the proof of the theorem.
\end{proof}

For the factorization of $f$, we assume the stably associated full
linear matrix $L$ is left monic. After we obtain atomic block
diagonal form as in Equation \ref{eq-block-atom}, we can factorize $L$
into atomic factors by Theorem~\ref{lfact1}. Combined with
Equation~\ref{eq4} we have
\begin{equation*}
f\oplus I_s = PS'F'_1F'_2\cdots F'_r U'Q,
\end{equation*}
where each linear matrix $F'_i$ is an atom, $P$ is upper triangular
with all $1$'s diagonal, $Q$ is lower triangular with all $1$'s
diagonal, and $S'$ is scalar invertible and $U'$ is a unit. Now,
applying Lemma \ref{extract} and Theorem \ref{cnzthm} we obtain the
complete factorization of $f$ into irreducible factors. This is
summarized in the following.

\begin{theorem}\label{comzthm}  
  Let $f\in \fR$ be a polynomial given by an algebraic branching
  program as input instance of $\prob{FACT}(\F_q)$. Then there is a
  $\poly(s, \log q,|X|)$ time randomized algorithm that outputs a
  factorization $f=f_1f_2\cdots f_r$ such that each $f_i$ is
  irreducible and is output as an algebraic branching program.
\end{theorem}

Analogous to Corollary \ref{sparse-cor1}, when the polynomial is given
in a sparse representation, we have 

\begin{corollary} \label{sparse-comnz}
If $f\in\fR$ is a polynomial given as input in sparse representation
(that is, an $\F_q$-linear combination of its monomials) then in
randomized polynomial time we can compute a factorization into
irreducible factors in sparse representation.
\end{corollary}

%
%
%
%
%
%
%
%
%
%
%
%
%
%
%
%
%
%
%



\subsection{Factorization over small finite fields}

Finally, we briefly discuss the factorization problem over small
finite fields.  As explained in Section \ref{small-field}, the two
steps in our factoring algorithm requiring randomization can be
replaced with deterministic $\poly(s, q, |X|)$ time
computation. Furthermore, as explained in Section~\ref{small-field},
the matrix shift $(M_1, M_2, \ldots, M_n)$ required for the Theorem
\ref{lfact2} can be obtained in deterministic polynomial time such
that the entries of the matrices $M_i$ are from $\F_q$ for each $i$.
Putting it together, it gives us a deterministic factorization
algorithm for noncommutative polynomials that are input as arithmetic
formulas over $\F_q$. In summary, we have the following.

\begin{theorem}
Given as input a multivariate polynomial $f\in\F_q\angle{X}$ for a
finite field $\F_q$ by a noncommutative algebraic branching program of size
$s$, a factorization of $f$ as a product $f=f_1f_2\cdots f_r$ can be
computed in deterministic time $\poly(s,q,|X|)$, where each
$f_i\in\F_q\angle{X}$ is an irreducible polynomial that is output as
an algebraic branching program.
\end{theorem}


